\section{Results}
\label{sec:results}

We report confirmed results first, then the \neurhlG{} gate evidence, then the
development-only season results and the null results. A result is
\emph{confirmed} only if it came from a window no decision had used, scored
once after the model was frozen. Results from windows that were consulted
during development are \emph{exploratory}.

\subsection{Confirmed: the player-game layer}
\label{sec:res-pg}

The player-game layer was frozen after development and scored once on the
player confirmation window, 2018--2020: 130,092 skater-games by 1,152 players in
3,624 games. The comparator is each skater's own recent history: an
exponentially weighted average ($\alpha = 0.1$) for ice-time share and shots,
renormalised over the dressed skaters, and for the two binary targets a
walk-forward shrunk average $\hat p = (n\,\bar x + k\,p_0)/(n + k)$ \citep{efron1975data}
with $k$ chosen on the previous season. This is a demanding baseline: in the first plan a neural
player-rate model had lost to it on ice time and shots (Table~\ref{tab:nulls}).
Table~\ref{tab:pg} shows the result. The layer is better on all four targets,
in every season, with two-way clustered $z$ statistics between $-8.1$ and
$-23.9$ and Holm-adjusted $p < 10^{-5}$.

\begin{table}[t]
\centering
\small
\caption{Player-game layer against each skater's own recent history on the
one-shot player confirmation, 2018--2020 (130,092 skater-games). Relative
difference in the stated loss; negative is better. $z$ uses standard errors
clustered by player and by game.}
\label{tab:pg}
\begin{tabular}{llrrc}
\toprule
Target & Loss & Relative difference & $z$ (two-way) & Seasons better \\
\midrule
Ice-time share & MAE & $-4.1\%$ & $-23.9$ & 3/3 \\
Shots on goal & Poisson deviance & $-10.2\%$ & $-13.7$ & 3/3 \\
P(goal $\geq 1$) & log loss & $-2.8\%$ & $-8.1$ & 3/3 \\
P(assist $\geq 1$) & log loss & $-2.6\%$ & $-15.9$ & 3/3 \\
\bottomrule
\end{tabular}
\end{table}

\subsection{Confirmed: \neurhlH{} against Elo}
\label{sec:res-h}

\neurhlH{} was frozen exactly as it had been evaluated on the tune window, the
freeze and the test were committed in an amendment before any 2018--2026 game
was scored by any \neurhlH{} artifact, and the test ran once. The primary set is
the eight structurally normal seasons of 2018--2026 ($n = 10{,}184$ games; 2021
excluded). Table~\ref{tab:c1} and Figure~\ref{fig:c1} give the results.

\begin{table}[t]
\centering
\small
\caption{One-shot confirmation of \neurhlH{} against Elo, regular-season games of
2018--2026 excluding 2021 ($n = 10{,}184$). Differences are \neurhlH{} minus Elo in
per-game log loss; negative favours \neurhlH{}.}
\label{tab:c1}
\begin{tabular}{lr}
\toprule
Quantity & Value \\
\midrule
Log loss, \neurhlH{} / Elo & 0.66495 / 0.66973 \\
Mean paired difference (SE) & $-0.00478$ (0.00088) \\
95\% CI, normal / moving-block bootstrap & \ci{-0.00650}{-0.00306} / \ci{-0.00660}{-0.00293} \\
Per-game test & $z = -5.46$, $p < 10^{-6}$ \\
Season-clustered $t$ (7 df) & $t = -5.71$, $p = 0.0007$ \\
Wild cluster bootstrap, exact ($2^8$ patterns) & $p = 0.0078$ \\
Seasons better; sign test & 8 of 8; $p = 0.0078$ \\
Including 2021 ($n = 11{,}052$) & $-0.00459$ \ci{-0.00622}{-0.00295}, 9 of 9 \\
\midrule
Brier score, \neurhlH{} / Elo & 0.23630 / 0.23857 \\
Reliability, \neurhlH{} / Elo & 0.00021 / 0.00025 \\
Resolution, \neurhlH{} / Elo & 0.01218 / 0.00982 \\
Uncertainty & 0.24840 \\
\bottomrule
\end{tabular}
\end{table}

\begin{figure}[t]
\centering
\includegraphics[width=\textwidth]{fig_c1_seasons}
\caption{\neurhlH{} minus Elo, per-game log loss by season on the one-shot
confirmation, with 95\% intervals; negative favours \neurhlH{}. Numbers under
the season labels are games. 2021 (hollow) is excluded from the primary test by
the declared rule and shown for completeness. The band and filled diamond are
the pooled primary estimate; the hollow diamond pools all nine seasons.}
\label{fig:c1}
\end{figure}

The effect is $-0.00478$ nats per game (95\% CI $-0.00650$ to $-0.00306$), and it
holds in every season. It is also the size the tune window had suggested
($-0.00283$), somewhat larger rather than smaller, and the amendment's power
table had put the chance of detecting an effect of the tune-window size at
about 0.93. Measured against a constant equal to the window's home-win rate
(0.540, log loss 0.68995), Elo's skill is 0.0202 nats and \neurhlH{}'s is 0.0250, an
increase of about 24\%. The Murphy decomposition locates the gain: reliability
is essentially unchanged and already small for both models, while resolution
rises from 0.0098 to 0.0122. \neurhlH{} is not a recalibrated Elo; it separates
games that Elo treats alike. Figure~\ref{fig:reliability} shows that both
models are well calibrated across the probability range.

The gain is conditional on information Elo does not use. \neurhlH{} reads which
skaters dressed for each game, which is public about an hour before puck drop,
and its Layer~1 translates those players' histories and career context into
projected shot shares. The confirmation therefore establishes that lineup-level
player information improves on a strong team-level rating system, not that a
network discovers team strength that Elo misses.

\paragraph{Where the gain comes from (exploratory).}
To test whether the lineup is what \neurhlH{} adds, we split the confirmation
games after the fact by how lopsided the two teams' absences were: the recent
ice time of regulars who did not dress, home minus away, in minutes. The window
is spent, so this analysis is exploratory. Its script stated two expectations
before it ran: that the gain would grow with the imbalance, and that it would
be largest early in the season, when rosters have changed and Elo has not yet
seen them play. The first held and the second did not (Table~\ref{tab:gain}).
Where neither team was missing regular ice time the difference is $-0.0003$
(SE 0.0045); it grows to $-0.0087$ when the imbalance exceeds 40 minutes, a
slope of $-0.0010$ per 10 minutes (SE 0.0004). By month the gain is smallest
in October and November and largest from December on. \neurhlH{} does not move
further from Elo when absences are lopsided (the mean absolute difference
between the two probabilities is 0.034 to 0.036 in every bin); its adjustments
are simply more accurate there.

\begin{table}[t]
\centering
\small
\caption{\neurhlH{} minus Elo, per-game log loss on the confirmation games,
split after the fact by the imbalance in absent regulars' recent ice time
between the two teams, and by month. Exploratory: the window had been spent.}
\label{tab:gain}
\begin{tabular}{lrr@{\hspace{2.5em}}lrr}
\toprule
Absence imbalance & Games & Difference (SE) & Month & Games & Difference (SE) \\
\midrule
none & 367 & $-0.0003$ (0.0045) & October & 1,316 & $-0.0017$ (0.0023) \\
0--20 min & 4,747 & $-0.0029$ (0.0013) & November & 1,726 & $-0.0022$ (0.0021) \\
20--40 min & 2,928 & $-0.0055$ (0.0016) & December & 1,674 & $-0.0078$ (0.0021) \\
over 40 min & 2,142 & $-0.0087$ (0.0020) & January & 1,600 & $-0.0057$ (0.0023) \\
 & & & February & 1,318 & $-0.0061$ (0.0024) \\
 & & & March & 1,712 & $-0.0042$ (0.0022) \\
 & & & April & 837 & $-0.0061$ (0.0032) \\
\bottomrule
\end{tabular}
\end{table}

\begin{figure}[t]
\centering
\includegraphics[width=0.62\textwidth]{fig_reliability}
\caption{Reliability of \neurhlH{} and Elo on the 10,184 confirmation games: ten
equal-count bins per model with 95\% Wilson intervals, the diagonal of perfect
calibration and the observed home-win rate (dashed); below, the distribution of
predicted home-win probabilities. \neurhlH{} spreads its forecasts more widely
while staying on the diagonal.}
\label{fig:reliability}
\end{figure}

\subsection{\neurhlG{}: iteration, gates and the pre-gate}
\label{sec:res-g}

\paragraph{Iteration.}
Table~\ref{tab:ladder} and Figure~\ref{fig:ladder} give the ladder on the
iteration window, every run of which is in the search ledger. The non-neural
engine (player-game sums through the integration and the stack) and the
tree-based rung did no better than Elo. The network with official power-play
targets (R6) improved on Elo by 0.0023 but did not reach \neurhlH{}. None of four
additions earned its place under the declared adoption rule (a gain of at
least 0.0003 without losing in four or more of six seasons): an Elo anchor on
the scoring rates, lineup-versus-usual multipliers, \neurhlH{}'s projections as
team inputs, and attention across the two rosters. Stacking R6 with \neurhlH{}
(R11) gave the best stack, 0.67258, still 0.00014 behind \neurhlH{} alone. The
paired standard error of each rung's difference from \neurhlH{} is 0.0007 to
0.0009, larger than most gaps between rungs, so the ladder shows the absence of
a large gain rather than a reliable ordering. The candidate declared for the
gate, g1, is R11 with the five-seed ensemble. The
candidate commit recorded, before the gate ran, the expectation that the
pre-gate would fail.

\begin{table}[t]
\centering
\small
\caption{\neurhlG{} iteration ladder on 2012 and 2014--2018 (7,421 games): pooled
log loss of each rung's final probability, from the search ledger. R4 was run
on the quick protocol (two seasons, 2,501 games), so its value is on a
different set of games. R6e and R6L were run before the master tensor was
rebuilt on official power-play opportunities; their parent on the earlier
tensor scored 0.67299, against which R6e is $+0.00031$ and R6L $-0.000004$, so
neither decision depends on the version.}
\label{tab:ladder}
\begin{tabular}{llr}
\toprule
Rung & Change & Log loss \\
\midrule
R1 & Elo & 0.67528 \\
R2 & \neurhlH{} & 0.67244 \\
R3 & gradient-boosted trees / logistic on game features, Elo-stacked & 0.67611 / 0.69699 \\
R4 & non-neural engine (quick protocol) & 0.67581 \\
R6 & network, official power-play targets & 0.67301 \\
R6e & + Elo anchor on scoring rates & 0.67330 \\
R6L & + lineup-versus-usual multipliers & 0.67299 \\
R6c & + \neurhlH{} projections as team inputs & 0.67344 \\
R13 & R6c + direct \neurhlH{} terms & 0.67344 \\
R9 & R6c + attention across rosters & 0.67366 \\
R11 & R6 stacked with \neurhlH{} (candidate structure) & 0.67258 \\
\bottomrule
\end{tabular}
\end{table}

\begin{figure}[t]
\centering
\includegraphics[width=\textwidth]{fig_g_ladder}
\caption{\neurhlG{} iteration ladder: pooled log loss of each full-run rung on the
iteration window (7,421 games), with Elo and \neurhlH{} on the same games. Open
circles were rejected under the adoption rule; the filled diamond, R11, was
carried forward as the candidate.}
\label{fig:ladder}
\end{figure}

\paragraph{Gates.}
The candidate was scored once on the gate window (6,289 games).
Table~\ref{tab:gates} summarises the gates and Figure~\ref{fig:gate} the
per-season differences. The final probability beats Elo by 0.00559 nats (SE
0.00131), in all five seasons, a margin comparable to \neurhlH{}'s confirmed one. It
beats \neurhlH{} by only 0.00049 (SE 0.00060), in three of five seasons. The
pre-gate required both a gain of at least 0.0045 over Elo and a gain of at
least 0.0010 over \neurhlH{}, each in at least four of five seasons. The first
condition holds and the second does not, so the pre-gate failed and, as the
plan declared, the sealed seasons were not spent. At the sealed test's size the
minimum detectable effect at 80\% power would have been 0.0057 against Elo and
0.0026 against \neurhlH{}.

\begin{table}[t]
\centering
\small
\caption{\neurhlG{} (candidate g1) on the gate window, 2019, 2020 and 2022--2024
(6,289 games; 225,752 skater-games). The calibration row reflects the
correction in Section~\ref{sec:failures}.}
\label{tab:gates}
\begin{tabular}{p{0.34\textwidth}p{0.56\textwidth}}
\toprule
Gate & Result \\
\midrule
Final probability vs Elo & $-0.00559$ (SE 0.00131), 5 of 5 seasons \\
Final probability vs \neurhlH{} & $-0.00049$ (SE 0.00060), 3 of 5 seasons \\
Pre-gate (S-STOP) & \textbf{fail}: the \neurhlH{} condition is not met \\
\midrule
PG: player heads vs the player-game layer & pass: ice-time share $-0.59\%$, shots $-0.45\%$ (Poisson log loss), P(goal) $-0.45\%$, P(assist) $-0.36\%$; every upper 95\% bound below zero (two-way clustered) \\
T: team box score & pass: shots, expected goals, goals and power-play opportunities beat log5 team history; shots and goals also beat the summed player-game layer, on deviance and CRPS \\
C: calibration & \textbf{fail}: slope 0.968 and overtime share (0.224 predicted, 0.220 observed) pass; randomised-PIT 80\% coverage passes for team goals (0.803) and skater shots (0.783) but not for team shots (0.860) \\
\bottomrule
\end{tabular}
\end{table}

\begin{figure}[t]
\centering
\includegraphics[width=\textwidth]{fig_g_gate_seasons}
\caption{\neurhlG{} on the gate window: per-season differences in per-game log loss
of the final probability against Elo (filled) and against \neurhlH{} (open), with
95\% intervals, and the pooled estimates over 6,289 games. Dotted segments mark
the pre-gate thresholds, $-0.0045$ against Elo and $-0.0010$ against \neurhlH{}.}
\label{fig:gate}
\end{figure}

\paragraph{The stat sheet.}
The engine's player heads improve on the confirmed player-game layer on all
four shared targets (Figure~\ref{fig:heads}). The gains are small, between
0.36\% and 0.59\%, but every upper confidence bound is below zero. Its team box
scores beat both declared baselines. Binned predictions track outcomes closely
for shots; for expected goals and goals they are somewhat too spread at the
extremes, so that the decile with the highest predicted regulation goals (4.2
per team-game) averaged 3.7 (Figure~\ref{fig:teambox}). The
win-probability slope is 0.968, and the hazard integration reproduces the
share of games that go past regulation without a tie scalar (22.4\% predicted
against 22.0\% observed; Figure~\ref{fig:otcal}). One calibration check fails:
the negative-binomial dispersion used for team shots makes those intervals too
wide, covering 86\% of outcomes at a nominal 80\%.

\begin{figure}[t]
\centering
\includegraphics[width=\textwidth]{fig_player_heads}
\caption{Relative change in loss of player-level predictions, with 95\%
intervals clustered by skater and game; negative is better, and the panels have
different scales. (a) The player-game layer against each skater's own history on
the one-shot confirmation, 2018--2020 (130,092 skater-games); shots on true
Poisson deviance. (b) \neurhlG{}'s player heads against the player-game layer on
the \neurhlG{} gate window (225,752 skater-games); shots on Poisson log loss. The
dotted line is the gate's non-inferiority margin of $+1\%$.}
\label{fig:heads}
\end{figure}

\begin{figure}[t]
\centering
\includegraphics[width=\textwidth]{fig_team_box}
\caption{\neurhlG{} team box-score predictions on the gate window (12,536 to
12,578 team-games): mean outcome by decile of prediction, with 95\% intervals, for
team shots on goal, expected goals and regulation goals, against the identity
line.}
\label{fig:teambox}
\end{figure}

\begin{figure}[t]
\centering
\includegraphics[width=\textwidth]{fig_outcome_calib}
\caption{\neurhlG{} outcome calibration on the gate window. (a) Reliability of
the final home-win probability in equal-count bins, with the calibration slope
and its 95\% interval. (b) Predicted and observed share of games tied after
regulation, by season and pooled; the prediction comes from the hazard
integration, with no tie scalar.}
\label{fig:otcal}
\end{figure}

\paragraph{Reading.}
Two independently built models that both condition on the dressed lineup, one
a small residual network over player histories feeding a four-feature logistic
head, the other a conservation-constrained engine that models every skater's
deployment and every team's scoring process, arrive at nearly the same
win-probability skill. On the iteration window the engine adds nothing beyond
\neurhlH{}, and on the gate window it adds an amount indistinguishable from zero.
We read this as evidence that the binding constraint is the information
available before the game, not the capacity or structure of the model; the
engine's value is that it produces a coherent box score at that same level of
skill.

\subsection{Development only: seasons and player seasons}
\label{sec:res-season}

The season simulator's 80\% intervals cover 0.794 of team-seasons, with a
standings error of 9.61 points (MAE) and a CRPS of 6.93, pooled over the
development and tune seasons. Its interval width was tuned on the same
seasons, so these are development numbers. On the five seasons that both its
record and the Elo baseline's share, the Elo baseline is more accurate: 9.17
against 9.99 points MAE and a CRPS of 6.43 against 7.17, ahead in four of five
seasons ($p = 0.19$). An earlier claim that the season layer beat the house
benchmarks compared different seasons and is withdrawn
(Section~\ref{sec:failures}). For player seasons, the declared 50/50 blend of a
gradient-boosted season model and the summed player-game layer had the lowest
points error over 11 vantages (9.07 points per 82 games, against 9.41 and 9.77
for the two paths alone).

\subsection{Documented nulls}
\label{sec:res-nulls}

Table~\ref{tab:nulls} lists the hypotheses that failed. Two deserve comment.
The first neural game model, a 1.2-million-parameter roster network over
player embeddings from an event language model, finished level with a constant
home-win rate. A regularised logistic regression on the same mean-pooled
embeddings did no better (0.6888 against 0.6898), while two team-form features
alone scored 0.6780, and adding the embeddings to form made it worse. A linear
model cannot be accused of failing to optimise, so the signal was absent from
the representation. The embeddings did learn something: position is linearly
decodable from them with 95\% accuracy, and linear probes recover
$R^2 = 0.54$ of points per 60 minutes and $0.59$ of ice time per game, although
neither was ever a training target. They encode what kind of player someone is, not how good his
team is, and every roster carries a similar mix of roles. The second is the
game-level reattempt of the third plan: four blocks of new information
(absence-based replacement value, goalie quality, schedule density, schedule
strength), each added to \neurhlH{}'s head in a declared order, each made the pooled
log loss worse, so no candidate reached the pre-gate.

\begin{table}[t]
\centering
\small
\caption{Documented null results. Each hypothesis was declared before its test
and failed on the stated evidence.}
\label{tab:nulls}
\begin{tabular}{p{0.44\textwidth}p{0.48\textwidth}}
\toprule
Hypothesis & Result \\
\midrule
A neural game network over event-model player embeddings beats Elo & Five configurations; best 0.6966 against Elo 0.67585 and a constant 0.6898 (tune window) \\
Pooled player embeddings carry team strength & Logistic probe 0.6888, level with the constant; embeddings added to team form make it worse \\
A career encoder can place rookies & Validation error 0.99--1.00 of predicting the mean at every vantage \\
Neural player-rate heads beat a player's own average & Worse on ice time (MAE 0.00934 against 0.00619) and shots; level with a shrunk average on goals \\
The event simulator's rates beat Elo per game & Blend 0.67763 against Elo 0.67668; its pre-gate stopped the confirmation \\
Absences, goalie quality, schedule density or strength improve the game head & Each block adds log loss (+0.00024 to +0.00541) \\
A starter model beats simple heuristics & Season start share wins (log loss 0.66--0.68 against 0.71--0.77) \\
Goals saved above expected predicts next season better than shrunk save percentage & $r = 0.159$ against $0.218$; Steiger $p = 0.007$ in favour of save percentage \\
RAPM transfers to a new team better than team-demeaned raw rates & Indistinguishable (movers $r = 0.387$ against $0.371$, $p = 0.60$) \\
The shot-quality model is calibrated to 0.005 per decile & Best 0.0065 after three declared correction ladders \\
\neurhlG{} adds win-probability information beyond \neurhlH{} & $-0.00049$ (SE 0.00060) on the gate window; the pre-gate failed \\
\bottomrule
\end{tabular}
\end{table}
